\documentclass[12pt,a4paper]{article}
\usepackage{cmap} % Makes the PDF copiable. See http://tex.stackexchange.com/a/64198/25761
\usepackage[T1]{fontenc}
\usepackage[brazil]{babel}
\usepackage[utf8]{inputenc}
\usepackage{amsmath}
\usepackage{amsfonts}
\usepackage{amssymb}
\usepackage{amsthm}
\usepackage{textcomp} % \degree
\usepackage{gensymb} % \degree
\usepackage[usenames,svgnames,dvipsnames]{xcolor}
\usepackage{hyperref}
\usepackage{graphicx}
\usepackage[margin=2cm]{geometry}

\hypersetup{
    colorlinks = true,
    allcolors = {blue}
}

% TODO: Consider using exsheets
% http://linorg.usp.br/CTAN/macros/latex/contrib/exsheets/exsheets_en.pdf
%
% http://ctan.org/tex-archive/macros/latex/contrib/exercise/
% Options: answerdelayed,lastexercise,noanswer
\usepackage[answerdelayed,lastexercise]{exercise}

\addto\captionsbrazil{%
\def\listexercisename{Lista de exerc\'icios}%
\def\ExerciseName{Exerc\'icio}%
\def\AnswerName{Solu\c{c}\~ao do exerc\'icio}%
\def\ExerciseListName{Ex.}%
\def\AnswerListName{Solu\c{c}\~ao}%
\def\ExePartName{Parte}%
\def\ArticleOf{de\ }%
}

\renewcommand{\ExerciseHeaderTitle}{(\ExerciseTitle)\ }
\renewcommand{\ExerciseListHeader}{%\ExerciseHeaderDifficulty%
\textbf{%\ExerciseListName\
\ExerciseHeaderNB.\ %
%\ --- \
\ExerciseHeaderTitle}%
%\ExerciseHeaderOrigin
\ignorespaces}
\renewcommand{\AnswerListHeader}{\textbf{\ExerciseHeaderNB.\ (\AnswerListName)\ }}

\newcommand*\sen{\operatorname{sen}}
\newcommand*\tg{\operatorname{tg}}
\newcommand*\dom[1]{\operatorname{Dom}\left(#1\right)}
\newcommand*\im[1]{\operatorname{Im}\left(#1\right)}

\newcommand*\R{\mathbb{R}}

\renewcommand{\theenumi}{\alph{enumi}}
\renewcommand\labelenumi{(\theenumi) }

\newcommand*\tipo{Prova I}
\newcommand*\turma{TADS121-01U}
\newcommand*\disciplina{CDI0001}
\newcommand*\eu{Helder G. G. de Lima}
\newcommand*\data{07/04/2018}

\author{\eu}
\title{\tipo - \disciplina}
\date{\data}

\begin{document}
\thispagestyle{empty}
\newgeometry{margin=2cm,bottom=0.5cm}
\begin{center}
\includegraphics[width=9.0cm]{marca} \\
\textbf{\tipo\ (\disciplina / \turma)} \\
Prof. \eu\footnote{
Este é um material de acesso livre distribuído sob os termos da licença \href{https://creativecommons.org/licenses/by-sa/4.0/deed.pt_BR}{Creative Commons Atribuição-CompartilhaIgual 4.0 Internacional}}
\end{center}

\noindent Nome do(a) aluno(a): \underline{\hspace{9,7cm}} Data: \underline{\data}

%\section*{Instruções}
\begin{center}\fbox{
\begin{minipage}{14cm}

{\footnotesize
\begin{itemize}
\renewcommand{\theenumi}{\Roman{enumi}}
\item Identifique-se em todas as folhas.
\item Mantenha o celular e os demais equipamentos eletrônicos desligados durante a prova.
\item Justifique cada resposta com cálculos ou argumentos baseados na teoria estudada.
\item Resolva (integralmente) apenas os itens de que precisar para somar 10,0 pontos.
\end{itemize}
}

\end{minipage}
}
\end{center}

%\section*{Questões}
\begin{ExerciseList}
\Exercise[title={2,0}] Seja $f(x) = \dfrac{1}{2-\sqrt[3]{x}}$. Obtenha a função $f^{-1}$, e exiba o domínio e a imagem de $f^{-1}$.
\Answer Considerando $y = f(x)$, tem-se $x = f^{-1}(y)$ e:
\begin{align*}
y = \dfrac{1}{2-\sqrt[3]{x}}
\Leftrightarrow y(2-\sqrt[3]{x}) = 1
\Leftrightarrow 2-\sqrt[3]{x} = \frac{1}{y}
%\Leftrightarrow -\sqrt[3]{x} = \frac{1}{y} - 2
\Leftrightarrow \sqrt[3]{x} = 2-\frac{1}{y}
\Leftrightarrow x = f^{-1}(y) = \left(2-\frac{1}{y}\right)^3.
\end{align*}
Como o denominador que aparece nas expressões acima não pode ser zero, conclui-se que:
\begin{itemize}
\item $\dom{f^{-1}} = \R^*$
\item $\im{f^{-1}} = \dom{f} = \R \setminus \{8\}$ (já que $2-\sqrt[3]{8} = 0$)
\end{itemize}
\Exercise[title={2,0}] Verifique se a função definida por $h(x) = (|x - 1| - |x + 1|)^2$ é par ou ímpar.
\Answer Observando que, para todo $x\in \R$, vale
\[
h(-x) = (|-x - 1| - |-x + 1|)^2
      = (|-(x + 1)| - |-(x - 1)|)^2
      = (|x + 1| - |x - 1|)^2
      = h(x),
\]
conclui-se que $h$ é par.
\Exercise[title={2,0}] O que é uma função periódica? Explique, dê um exemplo e confira se é injetora.
\Answer Uma função $f$ é periódica se existe $p \in \R^*$, tal que $f(x+p) = f(x)$, para todo $x \in \dom{f}$. Esse é o caso da função cosseno, já que $\cos(x + 2\pi) = \cos(x)$, para todo $x \in \R$. Como $\cos(0) = 1 = \cos(2\pi)$, e $0 \neq 2\pi$, conclui-se que a função cosseno não é injetora.
\Exercise[title={2,0}] Seja $q(x) = \dfrac{1}{\sqrt{9 - (x + 1)^2}}$.
\begin{enumerate}
\item Obtenha o intervalo $(a, b)$ formado por todos os pontos em que $q(x)$ está definida.
\item Determine $\displaystyle\lim_{x\to a^{+}} q(x)$.
\end{enumerate}

\Answer \begin{enumerate}
\item Para que $q(x)$ esteja definida, a expressão dentro da raiz quadrada não pode ser negativa, e o denominador não pode ser zero. Logo, basta identificar os valores de $x$ tais que $9 - (x + 1)^2 > 0$:
\begin{align*}
9 - (x + 1)^2 > 0
& \Leftrightarrow  9 > (x + 1)^2
  \Leftrightarrow \sqrt{9} > \sqrt{(x + 1)^2}
  \Leftrightarrow  3 > |x + 1|\\
& \Leftrightarrow -3 < x + 1 < 3
  \Leftrightarrow -4 < x < 2
  \Leftrightarrow x \in (-4, 2).
\end{align*}
\item Como $\lim_{x\to -4^{+}} 9 - (x + 1)^2 = 9 - (-4 + 1)^2 = 9 - (-3)^2 = 0$, e pelo item anterior $9 - (x + 1)^2 > 0$ para $-4 < x < 2$, então considerando $u = 9 - (x + 1)^2$, tem-se:
\[
  \lim_{x\to -4^{+}} \dfrac{1}{\sqrt{9 - (x + 1)^2}}
= \lim_{u\to 0^{+}} \dfrac{1}{\sqrt{u}}
= \lim_{v\to 0^{+}} \dfrac{1}{v}
= +\infty.
\]

\end{enumerate}
\Exercise[title={2,0}] Considerando que $r(x) = \dfrac{3x - 3}{x^2 - 1} + \dfrac{x^2 + 1}{x^3 + 1}$, calcule:
\begin{enumerate}
\item $\displaystyle \lim_{x \to +\infty} r(x)$
\item $\displaystyle \lim_{x \to 1} r(x)$
\end{enumerate}
\Answer
\begin{enumerate}
\item
\begin{align*}
\lim_{x \to +\infty} r(x)
& = \lim_{x \to +\infty} \dfrac{3x - 3}{x^2 - 1} + \dfrac{x^2 + 1}{x^3 + 1}
  = \lim_{x \to +\infty} \dfrac{3x - 3}{x^2 - 1}
+ \lim_{x \to +\infty} \dfrac{x^2 + 1}{x^3 + 1}\\
& = \lim_{x \to +\infty} \dfrac{3/x - 3/x^2}{1 - 1/x^2}
+ \lim_{x \to +\infty} \dfrac{1/x + 1/x^3}{1 + 1/x^3}\\
& = \dfrac{\lim_{x \to +\infty} 3/x - 3/x^2}{\lim_{x \to +\infty} 1 - 1/x^2}
+ \dfrac{\lim_{x \to +\infty} 1/x + 1/x^3}{\lim_{x \to +\infty} 1 + 1/x^3}\\
& = \dfrac{\lim_{x \to +\infty} 3/x - \lim_{x \to +\infty} 3/x^2}{\lim_{x \to +\infty} 1 - \lim_{x \to +\infty} 1/x^2}
+ \dfrac{\lim_{x \to +\infty} 1/x + \lim_{x \to +\infty} 1/x^3}{\lim_{x \to +\infty} 1 + \lim_{x \to +\infty} 1/x^3}\\
& = \dfrac{0 - 0}{1 - 0} + \dfrac{0 + 0}{1 + 0} = 0
\end{align*}
\item
\begin{align*}
\lim_{x \to 1} r(x)
& = \lim_{x \to 1} \dfrac{3x - 3}{x^2 - 1} + \dfrac{x^2 + 1}{x^3 + 1}
  = \lim_{x \to 1} \dfrac{3x - 3}{x^2 - 1} + \lim_{x \to 1} \dfrac{x^2 + 1}{x^3 + 1}
  = \lim_{x \to 1} \dfrac{3(x - 1)}{(x+1)(x-1) } + \lim_{x \to 1} \dfrac{x^2 + 1}{x^3 + 1}\\
&
  = \lim_{x \to 1} \dfrac{3}{(x+1) } + \lim_{x \to 1} \dfrac{x^2 + 1}{x^3 + 1}
  = \dfrac{3}{(1+1) } + \dfrac{1^2 + 1}{1^3 + 1} = \frac{5}{2}.
\end{align*}
\end{enumerate}

\Exercise[title={2,0}] Calcule os limites a seguir utilizando limites notáveis sempre que for possível:
\begin{enumerate}
\item $\displaystyle\lim_{x \to 0} \frac{\tg(2x)\sen(2x)}{x^2}$
\item $\displaystyle\lim_{x \to +\infty} \left(\frac{2x+3}{2x}\right)^{4x}$
\end{enumerate}
\Answer
\begin{enumerate}
\item
\begin{align*}
  \lim_{x \to 0} \frac{\tg(2x)\sen(2x)}{x^2}
& = \lim_{x \to 0} \tg(2x)\frac{\sen(2x)}{x^2}
  = \lim_{x \to 0} \frac{\sen(2x)}{\cos(2x)}\frac{\sen(2x)}{x^2}
  = \lim_{x \to 0} \frac{4}{\cos(2x)}\frac{\sen(2x)}{2x}\frac{\sen(2x)}{2x}\\
& = \lim_{u \to 0} \frac{4}{\cos(u)}\frac{\sen(u)}{u} \frac{\sen(u)}{u}
  = \lim_{u \to 0} \frac{4}{\cos(u)}
    \cdot\lim_{u \to 0} \frac{\sen(u)}{u}
    \cdot\lim_{u \to 0} \frac{\sen(u)}{u}\\
& = \frac{4}{\cos(0)}\cdot 1\cdot 1 = 4.
\end{align*}
\item
\begin{align*}
  \lim_{x \to +\infty} \left(\frac{2x+3}{2x}\right)^{4x}
& = \lim_{u \to +\infty} \left(\frac{u+3}{u}\right)^{2u}
= \lim_{u \to +\infty} \left(1+\frac{3}{u}\right)^{2u}
= \lim_{t \to +\infty} \left(1+\frac{1}{t}\right)^{6t}\\
& = \lim_{t \to +\infty} \left[\left(1+\frac{1}{t}\right)^{t}\right]^6
  = \lim_{w \to e} \left[w\right]^6 = e^6.
\end{align*}
\end{enumerate}

\Exercise[title={2,0}] Calcule os valores de $a$ e $b$ para que a seguinte função seja contínua:
\[
s(x) = \begin{cases}
\dfrac{\sqrt{x+1} - 1}{x}, & \text{se } x<0\\
a, & \text{se } x=0\\
\dfrac{b^x - 1}{x}, & \text{se } x > 0
\end{cases}
\]
\Answer Para que $s$ seja contínua em $0$, deve ocorrer $\lim_{x \to 0^-}s(x) = s(0) = \lim_{x\to 0^+} s(x)$. Pela direita, tem-se $\lim_{x \to 0^-} s(x) = \lim_{x \to 0^-} \dfrac{b^x - 1}{x} = \ln(b)$. Pela esquerda, tem-se:
\begin{align*}
\lim_{x \to 0^-} s(x)
& = \lim_{x \to 0^-} \dfrac{\sqrt{x+1} - 1}{x}
= \lim_{x \to 0^-} \dfrac{\sqrt{x+1} - 1}{x} \cdot \dfrac{\sqrt{x+1} + 1}{\sqrt{x+1} + 1}
  = \lim_{x \to 0^-} \dfrac{\sqrt{x+1}^2 - 1^2}{x \cdot (\sqrt{x+1} + 1)} \\
& = \lim_{x \to 0^-} \dfrac{x}{x \cdot (\sqrt{x+1} + 1)}
  = \lim_{x \to 0^-} \dfrac{1}{\sqrt{x+1} + 1}
  = \dfrac{1}{\sqrt{0+1} + 1}
  = \dfrac{1}{2}.
\end{align*}
Assim, $a = \ln(b) = 1/2$, isto é, $b = e^{1/2} = \sqrt{e}$.
\end{ExerciseList}

\begin{center}
BOA PROVA!
\end{center}

\newpage
\restoregeometry
\section*{Respostas}
\shipoutAnswer
\end{document}
